%
% File: abstract.tex
%
\chapter{Abstract}
\begin{SingleSpace}
The rapid maturation of generative artificial intelligence, and particularly of multimodal large language models capable of producing images, diagrams, and animated explanations on demand, has opened a new frontier in inclusive pedagogy. For learners with low language proficiency and for those with learning disabilities, the dependence of conventional instruction on dense text has long been a source of avoidable extraneous cognitive load. This systematic review synthesises ten recent empirical and theoretical studies, published between 2024 and 2025, that examine how AI-generated visual explanations function as a pedagogical resource for these populations. Drawing on the Cognitive Theory of Multimedia Learning and Cognitive Load Theory as analytic anchors, the review classifies the studies into four thematic clusters: AI tool typologies, theory--practice alignment, pedagogical efficacy, and implementation barriers. Findings indicate that AI-generated visual explanations are most effective when their design reflects established multimedia principles --- particularly the modality, signalling, and coherence principles --- and when they are embedded within scaffolded, learner-active tasks rather than offered as passive supplements. Evidence of efficacy is strongest for vocabulary learning, reading comprehension, and conceptual scaffolding in STEM, with reported gains of up to thirty per cent in some hybrid implementations. The review also identifies persistent gaps: a narrow disability focus on dyslexia and autism spectrum disorder, limited longitudinal data, and underdeveloped frameworks for managing AI hallucination and algorithmic bias in visual outputs. The paper concludes by proposing a research agenda that foregrounds adaptive load management, learner-controlled visual generation, and culturally responsive design.

\keywords{artificial intelligence; visual explanations; multimedia learning; cognitive load theory; learning disabilities; language proficiency; inclusive education; generative AI}
\end{SingleSpace}
% \clearpage
